%% SYNTHETIC TEST UNIT generated by writing/tools/make_outline_book.py from FINAL_OUTLINE_v2 -- not book text
\chapter{严格核心 I：$e^+e^-$ 湮灭总截面}\label{ch:ee}
本单元为测试用的合成文本（由 FINAL\_OUTLINE\_v2 自动生成），陈述内容从略。

\section{第 1 节（测试）}\label{sec:ee-R}

\begin{definition}[$R$ 比]\label{def:ee-R}
（测试用合成陈述：内容从略。）
% uses: none
\end{definition}

\begin{proposition}[$R=12\pi\,{\rm Im}\,\Pi$ \tagR]\label{prop:ee-R-ImPi}
（$R_{\mathrm{def}}$：由算符定义直接推出。）
（测试用合成陈述：内容从略。）
非微扰地位：\cref{asm:intro-wightman}。
\uses{def:ee-R, prop:prelim-optical, asm:intro-wightman}
\end{proposition}

\begin{lemma}[无质量传播子型图的单标度结构 \tagR]\label{lem:ee-single-scale}
（$R_{\mathrm{PT}}$：在重整化微扰论中对红外正规化的部分子态严格。）
（测试用合成陈述：内容从略。）
\uses{def:ir-landau-setup, def:qcd-dimreg}
\end{lemma}

\section{第 2 节（测试）}\label{sec:ee-analytic}

\begin{proposition}[解析性与正定性 \tagR]\label{prop:ee-analytic}
（$R_{\mathrm{def}}$：由算符定义直接推出。）
（测试用合成陈述：内容从略。）
非微扰地位：\cref{asm:intro-wightman}。
\uses{prop:ee-R-ImPi, asm:intro-wightman}
\end{proposition}

\begin{definition}[Adler 函数]\label{def:ee-adler}
（测试用合成陈述：内容从略。）
\uses{def:ee-R}
\end{definition}

\begin{proposition}[Adler 函数的色散表示 \tagR]\label{prop:ee-adler-dispersion}
（$R_{\mathrm{def}}$：由算符定义直接推出。）
（测试用合成陈述：内容从略。）
\uses{def:ee-adler, prop:ee-analytic}
\end{proposition}

\section{第 3 节（测试）}\label{sec:ee-euclid}

\begin{theorem}[欧氏红外有限性 \tagR]\label{thm:ee-euclid-ir}
（$R_{\mathrm{PT}}$：在重整化微扰论中对红外正规化的部分子态严格。）
（测试用合成陈述：内容从略。）
（本命题只涉及部分子态。）
\uses{thm:qcd-renormalizability, def:prelim-msbar, def:qcd-dimreg, prop:qcd-xi-independence}
\end{theorem}

\begin{proposition}[Adler 函数的重整化群 \tagR]\label{prop:ee-adler-rg}
（$R_{\mathrm{PT}}$：在重整化微扰论中对红外正规化的部分子态严格。）
（测试用合成陈述：内容从略。）
\uses{def:ee-adler, prop:prelim-rg, prop:qcd-rge}
\end{proposition}

\begin{theorem}[$R(s)$ 在微扰论每一阶的红外有限性 \tagR]\label{thm:ee-R-finite}
（$R_{\mathrm{PT}}$：在重整化微扰论中对红外正规化的部分子态严格。）
（测试用合成陈述：内容从略。）
非微扰地位：\cref{asm:intro-wightman}、\cref{asm:intro-bridge}、\cref{conj:ee-duality}。
\uses{prop:ee-R-ImPi, lem:ee-single-scale, thm:ee-euclid-ir, prop:ee-adler-rg}
\end{theorem}

\section{第 4 节（测试）}\label{sec:ee-theorem}

\begin{proposition}[$O(\alpha_s)$ 的实辐射与虚修正 \tagR]\label{prop:ee-oneloop}
（$R_{\mathrm{PT}}$：在重整化微扰论中对红外正规化的部分子态严格。）
（测试用合成陈述：内容从略。）
\uses{def:ee-R, prop:sudakov-oneloop, def:prelim-dr-conventions, lem:appC-phase-space}
\end{proposition}

\begin{example}[KLN 的作用与局限]\label{ex:ee-quark-spectrum}
（测试用合成陈述：内容从略。）
\uses{prop:ee-oneloop, thm:ir-kln}
\end{example}

\begin{proposition}[欧氏两圈关联函数 \tagR]\label{prop:ee-twoloop}
（$R_{\mathrm{PT}}$：在重整化微扰论中对红外正规化的部分子态严格。）
（测试用合成陈述：内容从略。）
\uses{def:prelim-msbar, lem:appC-bubble, lem:appC-ibp}
\end{proposition}

\section{第 5 节（测试）}\label{sec:ee-oneloop}

\begin{example}[欧氏两圈路线]\label{ex:ee-twoloop-route}
（测试用合成陈述：内容从略。）
\uses{prop:ee-twoloop, thm:ee-R-finite}
\end{example}

\begin{proposition}[$O(g^2)$ 处精确到 $\varepsilon$ 的切割恒等式 \tagR]\label{prop:ee-cutting-exact}
（$R_{\mathrm{PT}}$：在重整化微扰论中对红外正规化的部分子态严格。）
（测试用合成陈述：内容从略。）
\uses{prop:ee-oneloop, prop:ee-twoloop}
\end{proposition}

\begin{corollary}[与部分子末态求和的等同；KLN \tagA]\label{cor:ee-cuts}
（测试用合成陈述：内容从略。）
\uses{thm:ee-R-finite, thm:prelim-cutting-rules, rem:prelim-cutting-dr, prop:prelim-unitarity, prop:ee-cutting-exact}
\end{corollary}
\begin{proofsketch}测试。\end{proofsketch}
\paragraph{地位与缺口}测试：本陈述的地位由\cref{rem:prelim-cutting-dr}决定。

\section{第 6 节（测试）}\label{sec:ee-higher}

\begin{proposition}[围道公式 \tagR]\label{prop:ee-contour}
（$R_{\mathrm{def}}$：由算符定义直接推出。）
（测试用合成陈述：内容从略。）
\uses{prop:ee-adler-dispersion, prop:ee-analytic}
\end{proposition}

\begin{corollary}[$R(s)$ 的微扰展开 \tagR]\label{cor:ee-R-alphas}
（$R_{\mathrm{PT}}$：在重整化微扰论中对红外正规化的部分子态严格。）
（测试用合成陈述：内容从略。）
\uses{thm:ee-R-finite, prop:ee-oneloop, prop:ee-twoloop, prop:ee-adler-rg, prop:ee-contour}
\end{corollary}

\begin{example}[$O(a^3)$ 的 $\pi^2$ 项]\label{ex:ee-pi2}
（测试用合成陈述：内容从略。）
\uses{prop:ee-contour, cor:ee-R-alphas}
\end{example}

\section{第 7 节（测试）}\label{sec:ee-ope}

\begin{remark}[四个逻辑层次]\label{rem:ee-status-layout}
（测试用合成陈述：内容从略。）
% uses: none
\end{remark}

\begin{theorem}[流关联函数的欧氏 OPE \tagA]\label{thm:ee-ope}
（测试用合成陈述：内容从略。）
（本命题只涉及部分子态。）
\uses{thm:prelim-composite-renorm, thm:ee-euclid-ir, def:ee-adler}
\end{theorem}
\begin{proofsketch}测试。\end{proofsketch}
\paragraph{地位与缺口}测试：本陈述的地位由\cref{asm:intro-np-ope}、\cref{thm:dis-ope}决定。

\begin{conjecture}[非微扰 OPE 与凝聚 \tagC]\label{conj:ee-np-ope}
（测试用合成陈述：内容从略。）
\uses{thm:ee-ope}
\end{conjecture}
\paragraph{地位与缺口}测试：见文献。

\section{第 8 节（测试）}\label{sec:ee-duality}

\begin{remark}[复 $q^2$ 处的 OPE 是额外假设]\label{rem:ee-complex-ope}
（测试用合成陈述：内容从略。）
% uses: none
\end{remark}

\begin{proposition}[Poggio–Quinn–Weinberg 抹平恒等式 \tagR]\label{prop:ee-smearing}
（$R_{\mathrm{def}}$：由算符定义直接推出。）
（测试用合成陈述：内容从略。）
非微扰地位：\cref{asm:intro-wightman}。
\uses{prop:ee-analytic, asm:intro-wightman}
\end{proposition}

\begin{conjecture}[夸克—强子对偶 \tagC]\label{conj:ee-duality}
（测试用合成陈述：内容从略。）
\uses{prop:ee-smearing, conj:ee-np-ope}
\end{conjecture}
\paragraph{地位与缺口}测试：见文献。
